\documentclass[11pt]{letter}
\usepackage[margin=1in]{geometry}

\begin{document}

\begin{letter}{}

\opening{Dear Members of the Search Committee,}

I am writing to apply for the Assistant Professor position in Economics at the University of Virginia. I am a Ph.D. candidate in Economics at Johns Hopkins University. My research lies at the intersection of macroeconomics, public finance, and inequality, with a focus on using quantitative models and microdata to study distributional outcomes and policy design.

My job market paper develops and estimates a heterogeneous-agent model with differences in returns across households to study wealth inequality and the welfare effects of tax policy. Using micro-level data from the Survey of Consumer Finances, I estimate return distributions to match empirical wealth dynamics and simulate counterfactual policy environments. A key contribution of this work is showing how income risk and return heterogeneity jointly generate realistic wealth dispersion and marginal propensities to consume, allowing the model to match both inequality and aggregate behavior. More broadly, my research agenda combines structural modeling with micro and panel data, including work using the Health and Retirement Study where I apply econometric and causal identification strategies to study trust and financial returns.

In addition to my dissertation, I am developing collaborative work with a graduate researcher at UVA examining the distributional effects of oil price shocks on income, consumption inequality, employment, and asset prices. This project is motivated by the ``Innocent Bystanders'' literature and uses the Baumeister--Hamilton framework to identify oil supply and demand shocks. By linking macroeconomic shocks to micro-level outcomes, this work aims to better understand how global energy market fluctuations transmit through households and labor markets, with direct relevance for policy and inequality.

UVA would be an excellent environment to develop this research agenda. I see strong complementarities with faculty working in public finance, macroeconomics, and applied microeconomics, particularly Professor Lee Lockwood’s work on taxation and welfare, John Pepper’s work on inequality and empirical methods, and Eric Young’s work in quantitative macroeconomics. I would be especially excited to build on existing connections at UVA and contribute to collaborative work that integrates structural modeling, empirical analysis, and policy design.

I have teaching experience in Principles and Intermediate Macroeconomics, and I am comfortable teaching both undergraduate and graduate courses. My teaching emphasizes building economic intuition alongside quantitative tools, helping students connect theory to data and policy applications.

Thank you for your consideration. I would be thrilled to contribute to UVA’s research and teaching mission and to its intellectually vibrant community.

\closing{Sincerely,}

\vspace{0.5cm}

Decory Edwards

\end{letter}

\end{document}